\begin{abcliste}
\abc
Lastwagen und Auto drücken aufeinander: Die Kraft des Lastwagens auf das Auto und die Kraft des Autos auf den Lastwagen sind ein Wechselwirkungspaar. Nach dem dritten Newtonschen Gesetz sind sie
\[
 \result{\text{gleich gross und entgegengesetzt}}
\]
Das gilt immer, egal wie gross Massen und Geschwindigkeiten sind.

\abc
Tim hat nicht recht. Die Masse kommt im dritten Newtonschen Gesetz nicht vor: Der Lastwagen drückt genau so stark auf das Auto wie das Auto auf den Lastwagen. (Ebenso wenig drückt das schnellere Auto stärker.) Der schwerere Lastwagen ist nur schwerer abzubremsen: Dieselbe Kraft ändert seine Geschwindigkeit weniger. Darum \emph{wirkt} der Zusammenstoss auf das Auto viel stärker, obwohl die Kräfte gleich gross sind.

\abc
Zweites Newtonsches Gesetz, für jedes Fahrzeug mit seiner eigenen Masse und derselben Kraft $F$:
\begin{align}
 a_\mathrm{A} &= \aAF = \frac{\FK}{\mA} = \aA \approx \result{\aAP{0}{2}} \\
 a_\mathrm{L} &= \aLF = \frac{\FK}{\mL} = \aL \approx \result{\aLP{0}{2}}
\end{align}

\abc
Beide Kräfte wirken gleich lange ($\Delta t$), und $\Delta v = a\,\Delta t$. Die Geschwindigkeit des Autos ändert sich also stärker, und zwar
\begin{align}
 \frac{\Delta v_\mathrm{A}}{\Delta v_\mathrm{L}} &= \frac{a_\mathrm{A}}{a_\mathrm{L}} = \ratioF = \frac{\mL}{\mA} = \ratio \approx \result{\ratioP{0}{2}}
\end{align}
mal so stark wie die des Lastwagens.
\end{abcliste}
